\section{Optical implementation and discretization}
\label{sec:s:optics}

\paragraph{Coordinates and angular scan.}
The optical axis is $z$ and $\mathbf r=(x,y)$ is transverse position. In the simulated input medium, $n=1$ and $k_{\mathrm{ex}}=2\pi/\lambda_{\mathrm{ex}}$. The incident direction is
\begin{equation}
\mathbf k_t/k_{\mathrm{ex}}=
(\cos\theta_t\cos\alpha_t,\cos\theta_t\sin\alpha_t,\sin\theta_t).
\end{equation}
Thus, $\theta$ is the implemented vertical-plane parameter, not a polar angle from the optical axis. We use $\alpha_m=m\pi/16$, $m=0,\ldots,15$, with eight equally spaced $\theta$ values in each of $[\pi/6,\pi/3]$ and $[2\pi/3,5\pi/6]$. The second interval changes the sign of the transverse component while retaining positive longitudinal propagation. The resulting 256 directions are distinct. For $C=4$ and $8$, five fixed random subsets of 64 and 32 directions are shared across encoders.

\paragraph{Propagation and sampling.}
The incident field is uniformly $x$-polarized before phase modulation. Angular-spectrum propagation over $D=6.6\um$ carries the field from the plate to the intermediate illumination image plane. A finite NA-0.8 pupil filters the modulated angular spectrum before a vectorial Richards--Wolf relay generates three specimen-field components \cite{richards1959,marian2007}. Incident direction and transmitted angular spectrum are different quantities; the input direction range need not lie entirely inside the relay pupil. We sum the component intensities. Detection uses an NA-0.1 emission intensity PSF. Excitation and emission share $532$\,nm in this simulation; a physical implementation would use the actual spectral bands. No scattering, aberration or fluorescence saturation is modeled.

The specimen field is sampled laterally at $0.08\um$ and axially at $0.32\um$. Lateral binning gives an isotropic $0.32\um$ analysis grid of $32\times96\times96$ cells. The reconstruction basis averages $2\times2\times2$ analysis cells into each $0.64\um$ voxel. Camera pixels average $8\times8$ analysis samples laterally, giving $2.56\um$ pixel pitch and $12\times12$ samples per image. Distances and pitches are expressed in specimen-equivalent coordinates.

\paragraph{Image formation.}
Let $X_0(\mathbf r,z)$ denote fluorescence response per unit excitation and $h_{\mathrm{em}}$ the emission intensity PSF. With linear fluorescence and laterally shift-invariant detection, the noiseless camera-plane irradiance is
\begin{equation}
g_t(\mathbf r)=\int_{\Omega}h_{\mathrm{em}}(\mathbf r-\mathbf r',z_{\mathrm{focal}}-z')
X_0(\mathbf r',z')I_t(\mathbf r',z')\,\mathrm d^2\mathbf r'\,\mathrm dz'.
\label{eq:s:forward}
\end{equation}
Camera pixel $p$, occupying $D_p$, records $\bar y_{t,p}=|D_p|^{-1}\int_{D_p}g_t(\mathbf r)\,\mathrm d^2\mathbf r$. Emission adds incoherently over the specimen before this average. Excitation multiplies the specimen before detection blur; the detector need not resolve the illumination pattern itself.

\paragraph{Aperture and control encoders.}
The phase plate has $0.32\um$ features and initially independent phases uniform on $[0,2\pi)$. It is simulated on an aperture twice the specimen width, then the excitation is cropped to the specimen. Only the central plate region is optimized; its independent random surround remains fixed. The defocused amplitude controls use binary masks with the same feature pitch and 50\,\% open pixels at the same plate plane. The in-focus DMD uses independently drawn binary patterns relayed into focus. These are controls for the tested fixed-focus acquisition, not a comparison against all DMD microscopy configurations.

A finite illumination edge can itself encode depth when tilted. The plate-free control therefore uses an aperture four times the specimen width. Its residual central/border angular modulation is about 1.0/1.2\,\%, compared with 5.6/37\,\% at twice the width. The $4\times$-aperture plate-free study is kept separate from the $2\times$ phase/mask comparison. The random plate changes little across these apertures (centroid error $0.446$ versus $0.431\um$ at $C=8$, pretrained PnP), but that does not make their learned-versus-plate-free contrast a single-aperture experiment.

\paragraph{Integrated operator.}
The detection footprint $g_{p,i}$ includes the emission PSF, camera-pixel average and fine-cell integration weight. Reusing it in $A_{(t,p),j}=\sum_{i\in\mathcal P_j}I_t(\mathbf r_i,z_i)g_{p,i}$ avoids propagating each reconstruction basis element separately. Equality $\bA=\bB\bU$ is an algebraic statement for the selected piecewise-constant basis. It does not assert that a real specimen is constant within a voxel. For a fine-grid specimen $\bx_f$, the voxel-averaged reference is $\bx_{\mathrm{gt}}=\bR\bx_f$, where $\bR$ averages within reconstruction voxels. Normalized noisy measurements satisfy
\begin{equation}
\by=\bA\bx_{\mathrm{gt}}+
\underbrace{\bB(\bx_f-\bU\bR\bx_f)}_{\boldsymbol\delta_{\mathrm{disc}}}
+\boldsymbol\varepsilon.
\label{eq:s:patch}
\end{equation}
The retained residual is caused by specimen structure outside the reconstruction basis, not by an algebraic discrepancy between $\bA\bx$ and $\bB\bU\bx$.

\paragraph{Computational scope.}
After detection-footprint precomputation, direct assembly costs $O(TPF)$, compared with $O(TNF\log F)$ for constructing columns through separate FFT-based forward simulations: $T$ illumination patterns, $P$ camera samples, $F$ fine-grid cells and $N$ reconstruction voxels. Representative timings on an RTX 4500 Ada are approximately 20 minutes for 1000 optical-training steps, six minutes to assemble a full plate operator, 4--6 seconds for a PnP reconstruction at $C=8$, and 17 seconds for TV.
